% Options for packages loaded elsewhere
% Options for packages loaded elsewhere
\PassOptionsToPackage{unicode}{hyperref}
\PassOptionsToPackage{hyphens}{url}
\PassOptionsToPackage{dvipsnames,svgnames,x11names}{xcolor}
%
\documentclass[
  letterpaper,
  DIV=11,
  numbers=noendperiod]{scrreprt}
\usepackage{xcolor}
\usepackage{amsmath,amssymb}
\setcounter{secnumdepth}{5}
\usepackage{iftex}
\ifPDFTeX
  \usepackage[T1]{fontenc}
  \usepackage[utf8]{inputenc}
  \usepackage{textcomp} % provide euro and other symbols
\else % if luatex or xetex
  \usepackage{unicode-math} % this also loads fontspec
  \defaultfontfeatures{Scale=MatchLowercase}
  \defaultfontfeatures[\rmfamily]{Ligatures=TeX,Scale=1}
\fi
\usepackage{lmodern}
\ifPDFTeX\else
  % xetex/luatex font selection
\fi
% Use upquote if available, for straight quotes in verbatim environments
\IfFileExists{upquote.sty}{\usepackage{upquote}}{}
\IfFileExists{microtype.sty}{% use microtype if available
  \usepackage[]{microtype}
  \UseMicrotypeSet[protrusion]{basicmath} % disable protrusion for tt fonts
}{}
\makeatletter
\@ifundefined{KOMAClassName}{% if non-KOMA class
  \IfFileExists{parskip.sty}{%
    \usepackage{parskip}
  }{% else
    \setlength{\parindent}{0pt}
    \setlength{\parskip}{6pt plus 2pt minus 1pt}}
}{% if KOMA class
  \KOMAoptions{parskip=half}}
\makeatother
% Make \paragraph and \subparagraph free-standing
\makeatletter
\ifx\paragraph\undefined\else
  \let\oldparagraph\paragraph
  \renewcommand{\paragraph}{
    \@ifstar
      \xxxParagraphStar
      \xxxParagraphNoStar
  }
  \newcommand{\xxxParagraphStar}[1]{\oldparagraph*{#1}\mbox{}}
  \newcommand{\xxxParagraphNoStar}[1]{\oldparagraph{#1}\mbox{}}
\fi
\ifx\subparagraph\undefined\else
  \let\oldsubparagraph\subparagraph
  \renewcommand{\subparagraph}{
    \@ifstar
      \xxxSubParagraphStar
      \xxxSubParagraphNoStar
  }
  \newcommand{\xxxSubParagraphStar}[1]{\oldsubparagraph*{#1}\mbox{}}
  \newcommand{\xxxSubParagraphNoStar}[1]{\oldsubparagraph{#1}\mbox{}}
\fi
\makeatother


\usepackage{longtable,booktabs,array}
\newcounter{none} % for unnumbered tables
\usepackage{calc} % for calculating minipage widths
% Correct order of tables after \paragraph or \subparagraph
\usepackage{etoolbox}
\makeatletter
\patchcmd\longtable{\par}{\if@noskipsec\mbox{}\fi\par}{}{}
\makeatother
% Allow footnotes in longtable head/foot
\IfFileExists{footnotehyper.sty}{\usepackage{footnotehyper}}{\usepackage{footnote}}
\makesavenoteenv{longtable}
\usepackage{graphicx}
\makeatletter
\newsavebox\pandoc@box
\newcommand*\pandocbounded[1]{% scales image to fit in text height/width
  \sbox\pandoc@box{#1}%
  \Gscale@div\@tempa{\textheight}{\dimexpr\ht\pandoc@box+\dp\pandoc@box\relax}%
  \Gscale@div\@tempb{\linewidth}{\wd\pandoc@box}%
  \ifdim\@tempb\p@<\@tempa\p@\let\@tempa\@tempb\fi% select the smaller of both
  \ifdim\@tempa\p@<\p@\scalebox{\@tempa}{\usebox\pandoc@box}%
  \else\usebox{\pandoc@box}%
  \fi%
}
% Set default figure placement to htbp
\def\fps@figure{htbp}
\makeatother





\setlength{\emergencystretch}{3em} % prevent overfull lines

\providecommand{\tightlist}{%
  \setlength{\itemsep}{0pt}\setlength{\parskip}{0pt}}



 


\KOMAoption{captions}{tableheading}
\makeatletter
\@ifpackageloaded{bookmark}{}{\usepackage{bookmark}}
\makeatother
\makeatletter
\@ifpackageloaded{caption}{}{\usepackage{caption}}
\AtBeginDocument{%
\ifdefined\contentsname
  \renewcommand*\contentsname{Table of contents}
\else
  \newcommand\contentsname{Table of contents}
\fi
\ifdefined\listfigurename
  \renewcommand*\listfigurename{List of Figures}
\else
  \newcommand\listfigurename{List of Figures}
\fi
\ifdefined\listtablename
  \renewcommand*\listtablename{List of Tables}
\else
  \newcommand\listtablename{List of Tables}
\fi
\ifdefined\figurename
  \renewcommand*\figurename{Figure}
\else
  \newcommand\figurename{Figure}
\fi
\ifdefined\tablename
  \renewcommand*\tablename{Table}
\else
  \newcommand\tablename{Table}
\fi
}
\@ifpackageloaded{float}{}{\usepackage{float}}
\floatstyle{ruled}
\@ifundefined{c@chapter}{\newfloat{codelisting}{h}{lop}}{\newfloat{codelisting}{h}{lop}[chapter]}
\floatname{codelisting}{Listing}
\newcommand*\listoflistings{\listof{codelisting}{List of Listings}}
\makeatother
\makeatletter
\makeatother
\makeatletter
\@ifpackageloaded{caption}{}{\usepackage{caption}}
\@ifpackageloaded{subcaption}{}{\usepackage{subcaption}}
\makeatother
\usepackage{bookmark}
\IfFileExists{xurl.sty}{\usepackage{xurl}}{} % add URL line breaks if available
\urlstyle{same}
\hypersetup{
  pdftitle={Python Essentials with Hands-On Projects},
  pdfauthor={Shravan Parunandula},
  colorlinks=true,
  linkcolor={blue},
  filecolor={Maroon},
  citecolor={Blue},
  urlcolor={Blue},
  pdfcreator={LaTeX via pandoc}}


\title{Python Essentials with Hands-On Projects}
\usepackage{etoolbox}
\makeatletter
\providecommand{\subtitle}[1]{% add subtitle to \maketitle
  \apptocmd{\@title}{\par {\large #1 \par}}{}{}
}
\makeatother
\subtitle{A Top-Down Approach for Programming Aspirants}
\author{Shravan Parunandula}
\date{}
\begin{document}
\maketitle

\renewcommand*\contentsname{Table of contents}
{
\setcounter{tocdepth}{2}
\tableofcontents
}

\bookmarksetup{startatroot}

\chapter*{Preface}\label{preface}
\addcontentsline{toc}{chapter}{Preface}

\markboth{Preface}{Preface}

This is a Quarto book.

To learn more about Quarto books visit
\url{https://quarto.org/docs/books}.

\bookmarksetup{startatroot}

\chapter*{From Me to Every Programming
Aspirant}\label{from-me-to-every-programming-aspirant}
\addcontentsline{toc}{chapter}{From Me to Every Programming Aspirant}

\markboth{From Me to Every Programming Aspirant}{From Me to Every
Programming Aspirant}

I have mentored many students over the years, and I noticed one common
struggle almost everywhere:

Most students do not fail because programming is difficult. They
struggle because programming is often taught in a way that feels
disconnected from real understanding.

Many beginners are introduced to:

\begin{itemize}
\tightlist
\item
  syntax,
\item
  definitions,
\item
  theory,
\item
  and complex concepts
\end{itemize}

before they ever experience the joy of building something meaningful.

That is why I decided to create this handbook.

This is not just another Python notes collection. This handbook is
designed as a practical learning companion for every beginner who wants
to truly understand programming fundamentals.

I strongly believe students learn programming better when they:

\begin{itemize}
\tightlist
\item
  build first,
\item
  explore next,
\item
  and understand concepts naturally while solving problems.
\end{itemize}

So instead of starting with heavy theory, this handbook follows a
\textbf{Top-Down Approach}.

We begin with:

\begin{itemize}
\tightlist
\item
  small real-world projects,
\item
  simple applications,
\item
  hands-on coding,
\item
  and problem-solving exercises.
\end{itemize}

Then, step by step, we uncover:

\begin{itemize}
\tightlist
\item
  how the code works,
\item
  why each concept matters,
\item
  and how programmers think while building solutions.
\end{itemize}

My goal through this handbook is simple:

I want every student to:

\begin{itemize}
\tightlist
\item
  overcome the fear of coding,
\item
  think logically,
\item
  build confidence,
\item
  learn independently,
\item
  and enjoy the process of programming.
\end{itemize}

Inside this handbook, you will find:

\begin{itemize}
\tightlist
\item
  Python essentials explained simply,
\item
  beginner-friendly examples,
\item
  practical mini-projects,
\item
  coding exercises,
\item
  debugging mindset,
\item
  and real learning through implementation.
\end{itemize}

Programming is not about memorizing syntax.

Programming is about:

\begin{itemize}
\tightlist
\item
  thinking clearly,
\item
  solving problems,
\item
  experimenting,
\item
  making mistakes,
\item
  fixing them,
\item
  and growing through practice.
\end{itemize}

Every expert programmer once started as a beginner.

If you are someone who feels:

\begin{itemize}
\tightlist
\item
  ``Coding is hard,''
\item
  ``I cannot understand programming,''
\item
  or ``Maybe programming is not for me,''
\end{itemize}

then this handbook is especially for you.

Take one step at a time. Build small projects. Stay curious. Keep
practicing.

Because confidence in programming does not come from watching tutorials.

It comes from:

\begin{quote}
Building. Breaking. Fixing. Repeating.
\end{quote}

I hope this handbook helps you begin your journey with clarity,
confidence, and curiosity.

--- Shravankumar Parunandula

\part{Part I --- Getting Comfortable with Programming}

\chapter{Chapter 1 --- Welcome}\label{chapter-1-welcome}

\chapter*{Chapter 1 --- Welcome}\label{chapter-1-welcome-1}
\addcontentsline{toc}{chapter}{Chapter 1 --- Welcome}

\markboth{Chapter 1 --- Welcome}{Chapter 1 --- Welcome}

\emph{Placeholder content for Chapter 1 --- Welcome.}

\section{This is a placeholder for the content of Chapter 1. The actual
content will be added
later.}\label{this-is-a-placeholder-for-the-content-of-chapter-1.-the-actual-content-will-be-added-later.}

\subsection{This is a placeholder for the content of Chapter 1. The
actual content will be added
later.}\label{this-is-a-placeholder-for-the-content-of-chapter-1.-the-actual-content-will-be-added-later.-1}

\chapter{Chapter 2 --- Input and
Output}\label{chapter-2-input-and-output}

\chapter*{Chapter 2 --- Input and
Output}\label{chapter-2-input-and-output-1}
\addcontentsline{toc}{chapter}{Chapter 2 --- Input and Output}

\markboth{Chapter 2 --- Input and Output}{Chapter 2 --- Input and
Output}

\emph{Placeholder content for Chapter 2 --- Input and Output.}

\chapter{Chapter 3 --- Variables}\label{chapter-3-variables}

\chapter*{Chapter 3 --- Variables}\label{chapter-3-variables-1}
\addcontentsline{toc}{chapter}{Chapter 3 --- Variables}

\markboth{Chapter 3 --- Variables}{Chapter 3 --- Variables}

\emph{Placeholder content for Chapter 3 --- Variables.}

\chapter{Chapter 4 --- Conditions}\label{chapter-4-conditions}

\chapter*{Chapter 4 --- Conditions}\label{chapter-4-conditions-1}
\addcontentsline{toc}{chapter}{Chapter 4 --- Conditions}

\markboth{Chapter 4 --- Conditions}{Chapter 4 --- Conditions}

\emph{Placeholder content for Chapter 4 --- Conditions.}

\chapter{Chapter 5 --- Loops}\label{chapter-5-loops}

\chapter*{Chapter 5 --- Loops}\label{chapter-5-loops-1}
\addcontentsline{toc}{chapter}{Chapter 5 --- Loops}

\markboth{Chapter 5 --- Loops}{Chapter 5 --- Loops}

\emph{Placeholder content for Chapter 5 --- Loops.}

\part{Part II --- Thinking Like a Programmer}

\chapter{Chapter 6 --- Functions}\label{chapter-6-functions}

\chapter*{Chapter 6 --- Functions}\label{chapter-6-functions-1}
\addcontentsline{toc}{chapter}{Chapter 6 --- Functions}

\markboth{Chapter 6 --- Functions}{Chapter 6 --- Functions}

\emph{Placeholder content for Chapter 6 --- Functions.}

\chapter{Chapter 7 --- Lists}\label{chapter-7-lists}

\chapter*{Chapter 7 --- Lists}\label{chapter-7-lists-1}
\addcontentsline{toc}{chapter}{Chapter 7 --- Lists}

\markboth{Chapter 7 --- Lists}{Chapter 7 --- Lists}

\emph{Placeholder content for Chapter 7 --- Lists.}

\chapter{Chapter 8 --- Strings}\label{chapter-8-strings}

\chapter*{Chapter 8 --- Strings}\label{chapter-8-strings-1}
\addcontentsline{toc}{chapter}{Chapter 8 --- Strings}

\markboth{Chapter 8 --- Strings}{Chapter 8 --- Strings}

\emph{Placeholder content for Chapter 8 --- Strings.}

\chapter{Chapter 9 --- Dictionaries}\label{chapter-9-dictionaries}

\chapter*{Chapter 9 --- Dictionaries}\label{chapter-9-dictionaries-1}
\addcontentsline{toc}{chapter}{Chapter 9 --- Dictionaries}

\markboth{Chapter 9 --- Dictionaries}{Chapter 9 --- Dictionaries}

\emph{Placeholder content for Chapter 9 --- Dictionaries.}

\chapter{Chapter 10 --- Problem
Solving}\label{chapter-10-problem-solving}

\chapter*{Chapter 10 --- Problem
Solving}\label{chapter-10-problem-solving-1}
\addcontentsline{toc}{chapter}{Chapter 10 --- Problem Solving}

\markboth{Chapter 10 --- Problem Solving}{Chapter 10 --- Problem
Solving}

\emph{Placeholder content for Chapter 10 --- Problem Solving.}

\part{Part III --- Building Real Projects}

\chapter{Chapter 11 --- Quiz Game}\label{chapter-11-quiz-game}

\chapter*{Chapter 11 --- Quiz Game}\label{chapter-11-quiz-game-1}
\addcontentsline{toc}{chapter}{Chapter 11 --- Quiz Game}

\markboth{Chapter 11 --- Quiz Game}{Chapter 11 --- Quiz Game}

\emph{Placeholder content for Chapter 11 --- Quiz Game.}

\chapter{Chapter 12 --- Expense
Tracker}\label{chapter-12-expense-tracker}

\chapter*{Chapter 12 --- Expense
Tracker}\label{chapter-12-expense-tracker-1}
\addcontentsline{toc}{chapter}{Chapter 12 --- Expense Tracker}

\markboth{Chapter 12 --- Expense Tracker}{Chapter 12 --- Expense
Tracker}

\emph{Placeholder content for Chapter 12 --- Expense Tracker.}

\chapter{Chapter 13 --- Student Record
System}\label{chapter-13-student-record-system}

\chapter*{Chapter 13 --- Student Record
System}\label{chapter-13-student-record-system-1}
\addcontentsline{toc}{chapter}{Chapter 13 --- Student Record System}

\markboth{Chapter 13 --- Student Record System}{Chapter 13 --- Student
Record System}

\emph{Placeholder content for Chapter 13 --- Student Record System.}

\chapter{Chapter 14 --- Mini
Automation}\label{chapter-14-mini-automation}

\chapter*{Chapter 14 --- Mini
Automation}\label{chapter-14-mini-automation-1}
\addcontentsline{toc}{chapter}{Chapter 14 --- Mini Automation}

\markboth{Chapter 14 --- Mini Automation}{Chapter 14 --- Mini
Automation}

\emph{Placeholder content for Chapter 14 --- Mini Automation.}

\chapter{Chapter 15 --- Final Project}\label{chapter-15-final-project}

\chapter*{Chapter 15 --- Final
Project}\label{chapter-15-final-project-1}
\addcontentsline{toc}{chapter}{Chapter 15 --- Final Project}

\markboth{Chapter 15 --- Final Project}{Chapter 15 --- Final Project}

\emph{Placeholder content for Chapter 15 --- Final Project.}

\part{Part IV --- Beyond the Basics}

\chapter{Chapter 16 --- Files}\label{chapter-16-files}

\chapter*{Chapter 16 --- Files}\label{chapter-16-files-1}
\addcontentsline{toc}{chapter}{Chapter 16 --- Files}

\markboth{Chapter 16 --- Files}{Chapter 16 --- Files}

\emph{Placeholder content for Chapter 16 --- Files.}

\chapter{Chapter 17 --- Error Handling}\label{chapter-17-error-handling}

\chapter*{Chapter 17 --- Error
Handling}\label{chapter-17-error-handling-1}
\addcontentsline{toc}{chapter}{Chapter 17 --- Error Handling}

\markboth{Chapter 17 --- Error Handling}{Chapter 17 --- Error Handling}

\emph{Placeholder content for Chapter 17 --- Error Handling.}

\chapter{Chapter 18 --- Modules}\label{chapter-18-modules}

\chapter*{Chapter 18 --- Modules}\label{chapter-18-modules-1}
\addcontentsline{toc}{chapter}{Chapter 18 --- Modules}

\markboth{Chapter 18 --- Modules}{Chapter 18 --- Modules}

\emph{Placeholder content for Chapter 18 --- Modules.}

\chapter{Chapter 19 --- APIs}\label{chapter-19-apis}

\chapter*{Chapter 19 --- APIs}\label{chapter-19-apis-1}
\addcontentsline{toc}{chapter}{Chapter 19 --- APIs}

\markboth{Chapter 19 --- APIs}{Chapter 19 --- APIs}

\emph{Placeholder content for Chapter 19 --- APIs.}

\chapter{Chapter 20 --- Learning
Roadmap}\label{chapter-20-learning-roadmap}

\chapter*{Chapter 20 --- Learning
Roadmap}\label{chapter-20-learning-roadmap-1}
\addcontentsline{toc}{chapter}{Chapter 20 --- Learning Roadmap}

\markboth{Chapter 20 --- Learning Roadmap}{Chapter 20 --- Learning
Roadmap}

\emph{Placeholder content for Chapter 20 --- Learning Roadmap.}




\end{document}
